\chapter*{KATA PENGANTAR}
\addcontentsline{toc}{section}{
    \bfseries\protect\numberline{}KATA PENGANTAR
}

Puji dan syukur penulis panjatkan ke hadirat Tuhan Yang Maha Esa 
atas segala rahmat, berkat, dan karunia-Nya, sehingga Laporan Kemajuan 
Magang Industri ini dapat diselesaikan tepat pada waktunya. Laporan ini disusun sebagai salah satu syarat kelengkapan administrasi 
dan akademik dalam program Merdeka Belajar Kampus Merdeka (MBKM) 
Magang Industri bagi mahasiswa semester enam di Program Studi Teknik 
Informatika, Fakultas Teknologi Informasi, Universitas Tarumanagara. 

Pelaksanaan magang yang dilakukan di PT Inti Utama Solusindo 
(bagian dari Pharos Group) merupakan upaya strategis untuk menjembatani 
kompetensi teoritis yang diperoleh di bangku perkuliahan dengan 
kebutuhan praktis di industri teknologi informasi yang dinamis.

Penulis menyadari bahwa keberhasilan pelaksanaan magang dan 
penyusunan laporan ini tidak terlepas dari bimbingan, arahan, 
dan dukungan berbagai pihak. Oleh karena itu, penulis ingin 
menyampaikan apresiasi dan terima kasih yang sebesar-besarnya kepada:

\begin{enumerate}[label=\arabic*.,
    leftmargin=2.5em,
    labelindent=0pt,]
    \item Ibu Prof. Dr. Dyah Erny Herwindiati, M.Si., selaku Dekan Fakultas Teknologi Informasi Universitas Tarumanagara, yang telah memfasilitasi program MBKM bagi mahasiswa untuk berkembang di luar lingkungan kampus.
    \item Ibu Viny Christanti Mawardi, S.Kom., M.Kom., selaku Ketua Program Studi Teknik Informatika sekaligus Dosen Pembimbing, yang telah meluangkan waktu untuk memberikan arahan, saran, serta bimbingan yang sangat berharga selama proses magang berlangsung.
    \item Ibu Tony, S.Kom., M.Kom., Ph.D., selaku Dosen Koordinator Magang, yang telah memberikan informasi serta mengatur seluruh mekanisme administrasi magang industri di lingkungan Fakultas Teknologi Informasi.
    \item Saudara Dimas Aditya Nugroho, selaku mentor di PT Inti Utama Solusindo, yang telah memberikan kepercayaan, kesempatan, serta bimbingan teknis yang mendalam mengenai standar industri pengembangan perangkat lunak.
    \item Rekan-rekan mahasiswa satu grup bimbingan MBKM: Brandon Alexander Jayadi (Captain), Samuel Supardjo, Fico Filbert Budiman, Nicolas Matthew Wijaya, Mayceren, Monica Ho, Hans Thobie Sachio, Louis Febriano, Lusy Karlina Coa, Hengky Laurencio, Geraldus Raja Muka, dan Devany Aguslia, atas semangat, diskusi, dan kerja samanya.
\end{enumerate}

Penulis menyadari bahwa laporan kemajuan ini masih jauh dari kesempurnaan, 
baik dari segi kedalaman analisis maupun teknis penulisan. Terdapat 
berbagai tantangan yang masih harus dihadapi dalam sisa periode magang,
terutama dalam hal penguasaan kerangka kerja yang lebih kompleks. 
Namun, penulis berkomitmen untuk terus belajar dan memberikan kontribusi
terbaik bagi perusahaan. Semoga laporan ini dapat memberikan manfaat 
bagi para pembaca, khususnya bagi rekan-rekan mahasiswa yang akan menempuh 
program serupa di masa mendatang.

\begin{flushright}
    Jakarta, xx April 2026\\
    Peserta Magang Industri,\\[1.5cm]
    Aldian Yohanes\\
\end{flushright}